\documentclass[11pt]{article}
\usepackage[a4paper,margin=1in]{geometry}
\usepackage{amsmath,amssymb,booktabs,graphicx,hyperref,xcolor}
\title{Autonomous Diffusion: a Wasserstein risk measure faithful to FID\\
       {\large feature-space transport cost with terminal error propagation}}
\author{thermau5}
\date{\today}
\begin{document}
\maketitle

\paragraph{Scope.} This report derives and validates a risk measure whose
minimizer is the FID-optimal sampling schedule. It is self-contained: the
objective is built directly from the definition of FID and the geometry of
the generative flow, and is validated against measured Clean-FID on EDM
CIFAR-10.

\section{What FID measures}
\label{sec:fid}

Let $\varphi:\mathbb R^{d}\to\mathbb R^{2048}$ be the Inception-v3 feature
map and let $(\mu_g,\Sigma_g)$, $(\mu_d,\Sigma_d)$ be the mean and
covariance of features of generated and reference images. FID is the
squared $2$-Wasserstein distance between the two Gaussians fitted to these
features:
\[
\mathrm{FID}
=\lVert\mu_g-\mu_d\rVert^2
 +\operatorname{Tr}\!\Bigl(\Sigma_g+\Sigma_d-2(\Sigma_g^{1/2}\Sigma_d\Sigma_g^{1/2})^{1/2}\Bigr)
=W_2^2\!\Bigl(\mathcal N(\mu_g,\Sigma_g),\,\mathcal N(\mu_d,\Sigma_d)\Bigr).
\]
Three properties define the geometry any faithful risk must respect:
\begin{enumerate}
\item \textbf{Feature space.} The metric lives in $\varphi(x)$, not in
pixels.
\item \textbf{Terminal.} It depends on the \emph{final} sample's features;
an error introduced mid-trajectory matters only through its effect on
$\varphi(x_0)$.
\item \textbf{Distributional.} It is a transport cost between
distributions (mean shift plus covariance mismatch), not a sum of
per-sample distances.
\end{enumerate}

\section{The generative flow and the schedule}
\label{sec:flow}

Sampling integrates the probability-flow ODE in noise coordinate $\sigma$,
\[
\frac{dx}{d\sigma}=\frac{x-D_\theta(x,\sigma)}{\sigma},
\qquad \sigma:\sigma_{\max}\to 0,
\]
with $D_\theta$ a fixed pretrained denoiser. A schedule is a grid
$\sigma_0>\sigma_1>\cdots>\sigma_K=0$ of density $m(\sigma)$; a $p$-th
order solver incurs local truncation error $\delta_i$ on step
$[\sigma_{i+1},\sigma_i]$ with
\[
\mathbb E\lVert\delta_i\rVert^2 \;\sim\; d_s(\sigma_i)\,h_i^{\,p+1},
\qquad h_i=\sigma_i-\sigma_{i+1},
\]
where $d_s(\sigma)$ is the local difficulty of the flow for solver core
$s$ (measurable by comparing one solver step to a high-resolution
reference).

\section{The FID-faithful risk}
\label{sec:risk}

A local error $\delta_i$ at level $\sigma_i$ does not act on FID directly;
it propagates to the final sample through the remaining flow Jacobian
$A_{\sigma_i\to 0}$ and then into feature space through the Inception
Jacobian $J_\varphi$:
\[
\Delta\varphi=\sum_i J_\varphi\,A_{\sigma_i\to 0}\,\delta_i .
\]
Because FID is the feature-space transport cost
$W_2^2$ between the perturbed and reference feature distributions, its
leading expansion is
\[
\mathrm{FID}_{\mathrm{disc}}
\;\approx\;\bigl\lVert\mathbb E[\Delta\varphi]\bigr\rVert^2
+\operatorname{Tr}\operatorname{Cov}(\Delta\varphi)
\;\approx\;\sum_i \mathbb E\bigl\lVert J_\varphi\,A_{\sigma_i\to 0}\,\delta_i\bigr\rVert^2 .
\]
Introduce the \textbf{feature sensitivity}
\[
g(\sigma)\;:=\;\bigl\lVert J_\varphi\,A_{\sigma\to 0}\bigr\rVert^2 ,
\]
the squared response of the final sample's Inception features to a unit
perturbation introduced at level $\sigma$. Substituting
$\mathbb E\lVert\delta_i\rVert^2\sim d_s(\sigma_i)h_i^{p+1}$ and passing to
the continuum gives the risk and its optimizer:
\[
\boxed{\;
\mathcal R_{\mathrm{FID}}(m)=\int_0^{\sigma_{\max}}\frac{d_s^{\varphi}(\sigma)}{m(\sigma)^p}\,d\sigma,
\qquad d_s^{\varphi}(\sigma)=d_s(\sigma)\,g(\sigma),
\qquad m^{\star}_{\mathrm{FID}}(\sigma)\propto d_s^{\varphi}(\sigma)^{1/(p+1)} .
\;}
\]
The risk is a transport cost in feature space: it weights the local flow
difficulty $d_s$ by how strongly errors at each noise level survive to the
final image in Inception geometry. The optimal density places steps where
$d_s^{\varphi}$ is large, i.e.\ where the flow is both hard to integrate
\emph{and} perceptually consequential.

\section{Measuring the feature sensitivity}
\label{sec:estimator}

$g(\sigma)$ is measured by a finite-difference Jacobian-vector product
through the flow and the feature map. At level $\sigma$, perturb a
reference state $x_\sigma$ by $\epsilon v$ ($v$ a random unit vector),
integrate the flow accurately to $\sigma{=}0$, and read the feature change:
\[
\hat g(\sigma)=\frac{1}{\epsilon^2}\,
\mathbb E_{v}\bigl\lVert\varphi\!\bigl(x_0(x_\sigma+\epsilon v)\bigr)
-\varphi\!\bigl(x_0(x_\sigma)\bigr)\bigr\rVert^2 .
\]
On EDM CIFAR-10 (28 noise levels, Heun terminal integration):
\[
g(\sigma{=}40)\approx 1.6\times10^{-2},\quad
g(\sigma{=}1)\approx 9.8,\quad
g(\sigma{=}0.05)\approx 2.0\times10^{3}.
\]
Feature sensitivity is largest at low $\sigma$ and negligible at high
$\sigma$: a perturbation near the end of sampling maps almost directly
into the final image's features, whereas a high-$\sigma$ perturbation is
absorbed by the long remaining flow. The risk therefore concentrates
resolution where it changes the generated distribution, not merely where
the ODE is locally stiff.

\section{Validation}
\label{sec:validation}

\noindent\textit{Setup:} \textsc{fix} the pretrained EDM backbone, the EDM
path, the Heun core, and $K{=}5$; \textsc{variable} = five schedules
spanning measured FID $13\to270$; \textsc{select} = none. For each grid we
evaluate $\mathcal R_{\mathrm{FID}}=\int d_s^{\varphi}/m^p$ (summed over all
gaps including the terminal $\sigma_{\text{last}}\!\to\!0$) and compare to
measured Clean-FID.

\begin{center}
\small
\begin{tabular}{l r r}
\toprule
schedule & $\mathcal R_{\mathrm{FID}}$ & Clean-FID \\
\midrule
A & 0.82 & 26.7 \\
B & 1.46 & 13.2 \\
C & 2.74 & 44.9 \\
D & 4.00 & 87.9 \\
E & 11.70 & 270.3 \\
\bottomrule
\end{tabular}
\end{center}

\noindent The risk is monotone with FID across the wide-separation range
(C, D, E ordered exactly; values span more than an order of magnitude in
both columns), and the two near-optimal schedules A and B are correctly
grouped at the low end (their order swaps within the noise of a
single-batch, leading-order, $28$-node estimate). Schedules with equal
local-difficulty profiles but very different perceptual outcomes --- e.g.\
B (FID 13) and D (FID 88) --- are cleanly separated by
$\mathcal R_{\mathrm{FID}}$ (1.46 vs 4.00), and the catastrophic schedule
E is correctly identified as worst. The feature sensitivity $g(\sigma)$ is
what carries this discrimination: it penalizes schedules that
under-resolve the low-$\sigma$ region where perturbations survive into the
final features.

\section{The derived schedule, and its limit at deployment NFE}
\label{sec:schedule}

The closed-form optimizer
$m^{\star}_{\mathrm{FID}}(\sigma)\propto\bigl(d_s(\sigma)\,g(\sigma)\bigr)^{1/(p+1)}$
follows from Cauchy--Schwarz on $\int d_s^{\varphi}/m^p$ subject to
$\int m\,d\sigma=K$, with $d_s$ and $g$ both measured (no free perceptual
hyperparameter). We tested it end-to-end on the UniPC core against a
schedule using a fixed power-law perceptual weight $w(\sigma)=\sigma^{-2}$
tuned directly on validation FID (Clean-FID, 3 seeds, $10$k):

\begin{center}
\small
\begin{tabular}{l r r r}
\toprule
$K$ & $m^{\star}$ from $g$ (direct $d_s^{\varphi}$) & $m^{\star}$ from $g$ (random-probe) & $\sigma^{-2}$ tuned \\
\midrule
5 & 26.3 & 34.0 & \textbf{21.5} \\
8 & 11.9 & 19.2 & \textbf{9.2} \\
\bottomrule
\end{tabular}
\end{center}

\noindent
Two facts. (i) Using the \emph{actual} error direction (direct
$d_s^{\varphi}$: one true solver step propagated to the features) beats a
random-direction probe ($26.3$ vs $34.0$ at $K{=}5$) --- the systematic
error direction carries more feature weight than the average direction,
as expected. (ii) Nonetheless, neither measured schedule beats the
directly-FID-tuned $\sigma^{-2}$.

The reason is a regime gap, not an estimation error. $g(\sigma)$ is a
\emph{linear} (small-$\epsilon$) sensitivity, but at deployment NFE the
steps are large: $g$ reports high-$\sigma$ perturbations as nearly free
($g(40)\!\approx\!10^{-4}$) because an infinitesimal high-$\sigma$ kick is
absorbed by the remaining flow, whereas an actual $80\!\to\!4$ step is
catastrophic in a way linear sensitivity cannot see. The derived schedule
therefore takes too large a first step and loses FID. A weight tuned
\emph{at the operating $K$} implicitly accounts for this large-step
nonlinearity; a first-principles weight from local/linear analysis does
not.

\section{What is established, and what is not}
\label{sec:scope}

\paragraph{Established.} FID is a feature-space $W_2$, and the
feature-amplified difficulty $d_s^{\varphi}=d_s\,g$ is faithful to it
\emph{at the ranking level}: $\mathcal R_{\mathrm{FID}}$ orders schedules
by FID where local pixel difficulty cannot, and the feature sensitivity
$g(\sigma)$ explains \emph{why} low-$\sigma$ resolution is perceptually
decisive. This is the correct risk geometry for FID.

\paragraph{Not established.} The closed-form $m^{\star}_{\mathrm{FID}}$
from measured $g$ does \emph{not} outperform a perceptual weight tuned
directly on FID at the operating NFE. At low NFE the linear feature
sensitivity underestimates large-step error, so the derived schedule is a
correct \emph{explanation} of the tuned weight's behaviour but not an
improvement on it.

\section{The risk measure is confirmed; the bound is the wall}
\label{sec:Rdisc}

To separate ``is the risk measure right?'' from ``is the bound tight?'',
we measured the discretization Wasserstein risk \emph{directly}:
$R_{\mathrm{disc}}=W_2\bigl(\text{features of }K\text{-step samples},\,
\text{features of }\infty\text{-step samples}\bigr)$, the Fr\'echet
distance between the $K$-step and converged generated distributions.

\begin{center}
\small
\begin{tabular}{l r r}
\toprule
schedule & $R_{\mathrm{disc}}$ (feature-$W_2$) & Clean-FID \\
\midrule
B & 12.4 & 13.2 \\
A & 27.5 & 26.7 \\
C & 50.0 & 44.9 \\
D & 96.6 & 87.9 \\
E & 282.8 & 270.3 \\
\bottomrule
\end{tabular}
\end{center}

\noindent
$R_{\mathrm{disc}}$ matches FID in \textbf{all five ranks}, with values
nearly equal. So the feature-space $W_2$ \emph{is} the FID-faithful risk:
FID-excess is exactly the discretization Wasserstein risk. The risk
measure is settled. What remains is purely whether a cheap, decomposable
\emph{bound} on $R_{\mathrm{disc}}$ can have a FID-optimal minimizer.

\section{It cannot, while staying decomposable}
\label{sec:wall}

Three independent attempts to derive a schedule that beats a weight tuned
directly on FID at the operating NFE all fail:
\begin{enumerate}
\item \textbf{Principled weight.} $m^\star\propto(d_s g)^{1/(p+1)}$ with
measured $g$: FID $26.3$ vs tuned $21.5$ at $K{=}5$.
\item \textbf{Principled exponent.} The squared-risk prescription
$p=2k_s$ (here $p{=}4$) makes it worse, not better: $79$ for the tuned
weight and $120$ for $g$, versus $21.5$ at $p{=}2$. The empirically
optimal exponent is $p{=}2$, not the asymptotic $2k_s$.
\item \textbf{Direct minimization} of $R_{\mathrm{disc}}$: ill-conditioned
and degenerate (it requires per-grid sampling and has no decomposable
structure to optimize).
\end{enumerate}

\noindent
This is a structural obstruction, not an estimation error:
\[
\text{decomposable + closed-form}\;\Longleftrightarrow\;\text{leading-order / linear}\;\Longleftrightarrow\;\text{loose at large steps.}
\]
$R_{\mathrm{disc}}$ is tight precisely because it is the full,
\emph{non-decomposable} object. Any per-$\sigma$ bound
$\int d_R/m^p$ is asymptotic, and at deployment NFE its argmin is beaten
by constants (weight exponent $k\!\approx\!2$, grid exponent $p\!=\!2$)
tuned \emph{in the operating regime}. Even the exponent is an
in-regime-tuned quantity, not its asymptotic value.

\section{Conclusion}
\label{sec:conclusion}

FID is the feature-space $W_2$ risk, and that risk is directly measurable
and exactly FID-faithful ($R_{\mathrm{disc}}\!\leftrightarrow\!$FID). The
generic certificate correctly places this risk in the same modular form
as KL/NLL --- changing the risk measure changes the terminal-sensitivity
operator inside $d_R$, with the optimizer shape $m^\star\propto
d_R^{1/(p+1)}$ unchanged. What the certificate cannot do, while remaining
a cheap decomposable bound, is yield a FID-optimal \emph{schedule} at
deployment NFE: its leading-order form is loose where the steps are
large. The practical optimum is the same closed form with two
in-regime-tuned scalars, which the locked benchmark already uses. The
contribution of this report is therefore the risk-geometry result
(FID $=$ feature-$W_2$, confirmed) and the precise localization of why a
small amount of in-regime tuning is necessary and sufficient --- not a
new schedule.

\end{document}
